% It is suggested to use the "letterpaper" option below, output is
% currently unpredictable if anything else is chosen

% Including "titlepage" generates a separate titlepage and contents,
% leaving it out causes a title header to be placed on the first page

\documentclass[letterpaper, titlepage]{WARPdoc}

\begin{document}

\title{Using the WARPdoc Class}
\author{Gareth B Middleton}
\organization{Rice University\\ CMC Lab}

%This version number will be included in the footer on every page
%You must update it manually whenever you want a new version
% number to show up in the PDF
\version{1.2}

\maketitle

\section{Requirements}
The WARPdoc document class requires:
\begin{enumerate}

\item {\bf WARP logo}: The WARP project logo in PDF form must be in the root directory of your LaTeX project and must be named ``warplogo.pdf". A copy of the logo is included with this example document.

\item {\bf hyperref package}: The hyperref package must be in the path of your LaTeX installation. You can obtain a copy of the package from \href{http://www.tug.org/applications/hyperref/}{http://www.tug.org/applications/hyperref/}.

\item {\bf color \& graphicx packages}: These packages are required for the color formatting and graphics included in the WARPdoc class. These are standard packages which are probably included in your LaTeX installation.

\end{enumerate}

\section{Sections}
This document class supports hierarchy through sub-subsections, as illustrated below.

\subsection{Subsections}
This is the body of the first subsection.

\subsubsection{Sub-subsections}
This is the body of the first subsubsection.

\section{Figures}

Figures are included in the usual way. Here's a reference to the figure below: Figure \ref{fig:warplogo}.

\begin{figure}[h*]
	\centering
		\includegraphics{warplogo.pdf}
	\caption{This is the WARP logo}
	\label{fig:warplogo}
\end{figure}

\end{document}